\documentclass[11pt]{article}
\usepackage[margin=1in]{geometry}
\usepackage[T1]{fontenc}
\usepackage[utf8]{inputenc}
\usepackage{array}
\usepackage{longtable}
\usepackage{hyperref}
\usepackage{xcolor}
\usepackage{listings}
\lstset{basicstyle=\ttfamily\small,breaklines=true,columns=fullflexible,frame=single}
\setlength{\parindent}{0pt}
\setlength{\parskip}{0.65em}
\title{Lab 0: SimpleScalar Setup and Orientation}
\author{Computer Architecture Laboratory}
\date{}
\begin{document}
\maketitle
\section{Purpose}
This lab prepares students to build and run the SimpleScalar simulators used in the course. Students will verify the repository layout, build the simulators for PISA and Alpha, and run a small test program.
\section{Repository Layout}
After cloning the course repository, the repository root itself is the SimpleScalar working directory. If the repository is cloned as \texttt{SimpleScalar}, the layout is:
\begin{lstlisting}
SimpleScalar/
  benchmarks/
  config/
  labs/
  scripts/
  SimpleScalar/
    experiments/
  tests-alpha/
  tests-pisa/
  sim-fast.c
  sim-safe.c
  sim-profile.c
  sim-bpred.c
  sim-outorder.c
  sim-singlecycle.c
  Makefile
\end{lstlisting}
There is no extra nested \texttt{simplesim-3.0/} directory in this repository. Run the commands in these manuals from the cloned repository root unless a lab explicitly says to enter \texttt{labs/}.
\section{Simulator Overview}
\begin{center}
\scriptsize
\begin{longtable}{|p{0.24\linewidth}|p{0.24\linewidth}|}
\hline
Simulator & Purpose \\
\hline
\endfirsthead
\hline
Simulator & Purpose \\
\hline
\endhead
\texttt{sim-fast} & Fast functional simulation with minimal checking. \\
\hline
\texttt{sim-safe} & Functional simulation with additional safety checks. \\
\hline
\texttt{sim-profile} & Instruction and address profiling. \\
\hline
\texttt{sim-cache} & Cache hierarchy simulation. \\
\hline
\texttt{sim-bpred} & Branch predictor simulation. \\
\hline
\texttt{sim-outorder} & Detailed out-of-order processor timing simulation. \\
\hline
\texttt{sim-singlecycle} & Educational single-cycle timing model added for this course. \\
\hline
\end{longtable}
\normalsize
\end{center}
\section{Part A: Build for PISA}
Open a terminal in the cloned repository root:
\begin{lstlisting}
cd /path/to/SimpleScalar
\end{lstlisting}
Clean any old build files:
\begin{lstlisting}
make clean
\end{lstlisting}
Configure for PISA:
\begin{lstlisting}
make config-pisa
\end{lstlisting}
Compile the simulators:
\begin{lstlisting}
make
\end{lstlisting}
Check that simulator executables were created:
\begin{lstlisting}
ls sim-fast sim-safe sim-profile sim-bpred sim-outorder sim-singlecycle
\end{lstlisting}
Run a quick PISA test:
\begin{lstlisting}
./sim-safe tests-pisa/bin.little/test-math
\end{lstlisting}
\section{Part B: Build for Alpha}
Clean the previous build:
\begin{lstlisting}
make clean
\end{lstlisting}
Configure for Alpha:
\begin{lstlisting}
make config-alpha
\end{lstlisting}
Compile again:
\begin{lstlisting}
make
\end{lstlisting}
Run a quick Alpha test:
\begin{lstlisting}
./sim-safe tests-alpha/bin/test-math
\end{lstlisting}
\section{Part C: Getting Help From a Simulator}
Every simulator supports a help option:
\begin{lstlisting}
./sim-profile -h
./sim-outorder -h
./sim-bpred -h
\end{lstlisting}
Use the help output to confirm the option names used in later labs.
\end{document}
